\chapter{Diferenciación}
\label{cap:diferenciacion}

La derivada se define mediante un límite y se estudia qué información local aporta. Los teoremas de Rolle y del valor medio fundamentan los principales criterios y aplicaciones de la diferenciación \cite{spivak2008,apostol1974,rudin1976}.

\section{Definición de derivada}
Cociente incremental, derivada en un punto y derivabilidad en un intervalo.

\section{Interpretación local de la derivada}
Aproximación lineal, tangente y relación entre derivabilidad y continuidad.

\section{Reglas de diferenciación}
Reglas de suma, producto y cociente, con demostraciones.

\section{Regla de la cadena y derivadas de funciones inversas}
Regla de la cadena y fórmula para la derivada de una función inversa.

\section{Derivadas de orden superior}
Derivadas iteradas y notación para órdenes superiores.

\section{Teorema de Rolle}
Existencia de un punto crítico bajo las hipótesis de Rolle.

\section{Teorema del valor medio}
Enunciado, prueba y aplicaciones a monotonía y estimaciones.

\section{Teorema del valor medio de Cauchy}
Versión generalizada y relación con el teorema de Rolle.

\section{Consecuencias del teorema del valor medio}
Unicidad, acotación de incrementos y criterios de constancia o monotonía.

\section{Regla de L'Hôpital}
Resolución de formas indeterminadas mediante cocientes de derivadas bajo las hipótesis adecuadas.

\section{Funciones trascendentes}
Construcción y propiedades de las funciones logarítmica, exponencial, trigonométricas e inversas.

\subsection{La función logarítmica}
Definición y propiedades derivadas de la integral de $1/x$.

\subsection{La función exponencial}
Definición, propiedades algebraicas y relación con la función logarítmica.

\subsection{Funciones trigonométricas}
Definiciones y derivadas de seno, coseno y las demás funciones trigonométricas.

\subsection{Funciones trigonométricas inversas}
Dominios, rangos y derivadas de las funciones trigonométricas inversas.

\section{Concavidad de funciones}
Uso de la segunda derivada para estudiar convexidad, concavidad y puntos de inflexión.

\section{Fórmula y teorema de Taylor}
Polinomios de Taylor y representación del error en la aproximación local de una función.

\section{Formas del resto y estimaciones del error}
Resto de Lagrange, forma integral y cotas del error con aplicaciones cuantitativas.
